\section{Actual formal initial diagrams and their finite monic vertex series}

Let $S=\mathbb C[[X_1,\ldots,X_n]]$ be the GENUINE formal series ring.
Fix an additive, pointwise-monotone well monomial order $\preceq$ on the
actual exponent set $\mathbb N^n$. These are actual coefficient supports
and the actual finite-antidiagonal product. This is the initial-diagram
construction used in \cite[\S3, \S7, pp.~261--264]{bierstone-milman97-invent-math},
before proving formal division. The degree-lexicographic specialization
uses a fixed ordering of the original coordinates; reversing the coordinate
priority only permutes the convention permitted in the source.

\begin{lemma}[Actual lowest nonzero coefficients and genuine convolution]
\label{hir:actual-formal-initial-exponent-coefficient-product}
For every nonzero actual $F\in S$, there is a unique lowest exponent
$\exp_\preceq(F)$ in its actual nonzero coefficient support. Its actual
coefficient is nonzero and every smaller coefficient vanishes. For actual
nonzero $F,G$,
\[
 \exp_\preceq(FG)=\exp_\preceq(F)+\exp_\preceq(G),\qquad
 [X^{\exp(F)+\exp(G)}](FG)
 =[X^{\exp(F)}]F\,[X^{\exp(G)}]G.
\]
In particular the lowest exponent of the genuine monomial $X^\beta$ is
$\beta$, and multiplying $F$ by it shifts its actual lowest exponent by
$\beta$.
\lean{Hironaka.HironakaActualFormalInitialDiagram.exists_actualInitialExponent,Hironaka.HironakaActualFormalInitialDiagram.actualInitialExponent,Hironaka.HironakaActualFormalInitialDiagram.actualInitialExponent_spec,Hironaka.HironakaActualFormalInitialDiagram.coeff_eq_zero_of_lt_actualInitialExponent,Hironaka.HironakaActualFormalInitialDiagram.actualInitialExponent_eq_of_coefficients,Hironaka.HironakaActualFormalInitialDiagram.coeff_mul_actualInitialExponent,Hironaka.HironakaActualFormalInitialDiagram.actualInitialExponent_mul,Hironaka.HironakaActualFormalInitialDiagram.actualMonomial_ne_zero,Hironaka.HironakaActualFormalInitialDiagram.actualInitialExponent_monomial}
\leanok
\end{lemma}
\begin{proof}
The actual nonzero coefficient support is nonempty, so its image in the
genuine well ordered exponent type has a minimum. Pull it back to the
same exponent set. Membership gives its nonzero coefficient; minimality
gives vanishing below it. Conversely these two properties characterize
that exponent by antisymmetry.
Put $a=\exp(F)$ and $b=\exp(G)$. In the ACTUAL finite convolution
for degree $a+b$, the pair $(a,b)$ contributes the displayed product.
For every other pair $(u,v)$ with $u+v=a+b$, either $u\prec a$ or
$v\prec b$: if both were at least their respective minima, ordered
addition and cancellation force equality of both. Thus every other
actual summand vanishes. The displayed product is nonzero over
$\mathbb C$. If $c\prec a+b$, then for every pair $u+v=c$ one of
$u\prec a$, $v\prec b$ again holds by ordered addition; hence that
whole convolution vanishes. The characterization proves the product
formula. The monomial has precisely its one actual nonzero coefficient,
so the same characterization and product formula prove the final claim.
\end{proof}

\begin{theorem}[The whole genuine ideal determines its actual upper diagram]
\label{hir:actual-formal-ideal-initial-diagram-upper}
\uses{hir:actual-formal-initial-exponent-coefficient-product}
For every actual ideal $I\subset S$, define its actual initial diagram
\[
 \mathcal N_\preceq(I)=
 \{\exp_\preceq(F):F\in I,\ F\ne0\}.
\]
Then $a\in\mathcal N_\preceq(I)$ implies
$a+\beta\in\mathcal N_\preceq(I)$ for every actual exponent $\beta$.
Equivalently, this actual diagram is upward closed for the original
pointwise exponent order. The zero ideal has empty diagram.
\lean{Hironaka.HironakaActualFormalInitialDiagram.actualInitialDiagram,Hironaka.HironakaActualFormalInitialDiagram.actualInitialDiagram_add,Hironaka.HironakaActualFormalInitialDiagram.actualInitialDiagram_upper,Hironaka.HironakaActualFormalInitialDiagram.actualInitialDiagram_bot}
\leanok
\end{theorem}
\begin{proof}
Take the actual ideal member $F$ witnessing $a$. The actual monomial
$X^\beta$ times $F$ is a member of the SAME whole ideal, is nonzero,
and has lowest exponent $\beta+a$ by the preceding result. If $a\le c$
pointwise, write $c=(c-a)+a$ in the actual finitely supported exponent
monoid and apply this assertion. The zero ideal has no actual nonzero
member, proving its empty-diagram statement directly.
\end{proof}

\begin{theorem}[Dickson vertices have internally constructed monic genuine series]
\label{hir:actual-formal-initial-diagram-finite-monic-vertices}
\uses{hir:actual-formal-ideal-initial-diagram-upper,hir:actual-formal-initial-exponent-coefficient-product}
For every actual $I\subset S$, its pointwise-minimal diagram elements
form a finite set $V$. For EVERY actual exponent $c$,
\[
 c\in\mathcal N_\preceq(I)\quad\Longleftrightarrow\quad
 \exists a\in V:\ a\le c.
\]
For each $a\in V$, one can internally choose a GENUINE $f_a\in I$
whose actual lowest exponent is $a$ and whose actual coefficient there
is $1$. In particular all these $f_a$ are nonzero. No vertex set or
series with the required exponents is a premise.
\lean{Hironaka.HironakaActualFormalInitialDiagram.actualInitialDiagram_vertices_finite,Hironaka.HironakaActualFormalInitialDiagram.actualInitialVertices,Hironaka.HironakaActualFormalInitialDiagram.mem_actualInitialDiagram_iff_vertex_le,Hironaka.HironakaActualFormalInitialDiagram.exists_actual_monic_vertex_series}
\leanok
\end{theorem}
\begin{proof}
Dickson's lemma applies to the SAME finite-coordinate pointwise
exponent order. Minimal elements form an antichain, hence are finite.
For each actual diagram member $c$, pointwise well foundedness gives
a minimal diagram member $a\le c$; conversely upper closure of the
actual diagram proves membership from any such $a$.
For each $a$ in this finite set, actual diagram membership supplies an
actual nonzero $F_a\in I$ with lowest exponent $a$. Its coefficient
$c_a=[X^a]F_a$ is a genuine nonzero complex number. Set
$f_a=C(c_a^{-1})F_a$ in the ACTUAL formal series ring. The whole
ideal is closed under this actual product. The true coefficient formula
for a constant series product makes its coefficient at $a$ equal to
$1$ and preserves vanishing of every lower coefficient. The preceding
characterization therefore gives the SAME lowest exponent $a$.
This constructs every claimed series internally. It does not yet claim
that this family generates the whole ideal: that is a consequence of
the subsequent genuine formal division theorem.
\end{proof}

\begin{lemma}[Degree-compatible lowest exponents detect actual formal powers]
\label{hir:actual-formal-degreelex-initial-exponent-power-kernel}
\uses{hir:actual-formal-initial-exponent-coefficient-product,hir:actual-formal-variable-power-coefficient-kernel}
Specialize to the genuine degree-lexicographic monomial order. For each
actual nonzero $F\in S$ and each $q\ge0$,
\[
 F\in(X_1,\ldots,X_n)^q\quad\Longleftrightarrow\quad
 |\exp_{\mathrm{deglex}}(F)|\ge q.
\]
\lean{Hironaka.HironakaActualFormalInitialDiagram.mem_formal_variables_pow_iff_degree_initialExponent}
\leanok
\end{lemma}
\begin{proof}
If $F$ lies in the true variable power, the already proved actual power
kernel kills every coefficient below total degree $q$. Its genuine
lowest coefficient is nonzero, so its degree is at least $q$.
Conversely suppose that degree is at least $q$. Every exponent of
degree below $q$ is strictly smaller in the genuine degree-first order,
so its actual coefficient vanishes by the lowest-coefficient lemma.
The BOTH-direction genuine power kernel places $F$ in that same actual
formal variable power. This includes $q=0$ and the empty coordinate
family; it assumes no Hilbert--Samuel or standard-basis certificate.
\end{proof}
